\section*{Capítulo \ref{ch:g7:sound-production} --- Som: produção e propagação}

\begin{solution}{exo:g7:sound-production:1}
É a substância pela qual um \omterm{def:g3:sound-around-us:sound}{som} viaja --- \omterm{def:g2:air-around-us:air}{ar}, água, madeira, aço,
qualquer multidão molecular. O não \omterm{def:g7:sound-production:medium}{meio}: o vácuo.
\end{solution}

\begin{solution}{exo:g7:sound-production:2}
Quem viaja é o revezamento: um padrão de apertos em disparada,
passado de fileira em fileira da multidão. Cada \omterm{def:g7:states-of-matter:molecule}{molécula} apenas se
sacode em torno da casa dela --- a mensagem se move, os mensageiros
ficam.
\end{solution}

\begin{solution}{exo:g7:sound-production:3}
Sem multidão, sem revezamento: no frasco com o \omterm{def:g2:air-around-us:air}{ar} bombeado para fora,
as paredes trêmulas da \omterm{ex:g6:circuits-and-safety:motor}{campainha} empurram o nada, e mensagem nenhuma
se forma. As batalhas espaciais são silenciosas pelo mesmo motivo ---
explosões entre as \omterm{def:g4:solar-system:star}{estrelas} não têm multidão para carregar o estrondo
delas.
\end{solution}

\begin{solution}{exo:g7:sound-production:4}
\omterm{def:g2:air-around-us:air}{Ar}, cerca de \qty{340}{m/s}; água, cerca de \qty{1500}{m/s}; aço,
cerca de \qty{5000}{m/s}. Quanto mais apertadas se tocam as
\omterm{def:g7:states-of-matter:molecule}{moléculas}, mais depressa o empurrão é passado adiante: voando,
encostadas, travadas.
\end{solution}

\begin{solution}{exo:g7:sound-production:5}
As fileiras travadas do aço revezam quase de mão em mão: cerca de
quinze vezes o ritmo do \omterm{def:g2:air-around-us:air}{ar}, e por isso a mensagem do trilho chega
muito antes do ronco pelo \omterm{def:g2:air-around-us:air}{ar}.
\end{solution}

\begin{solution}{exo:g7:sound-production:6}
A viagem da \omterm{def:g1:light-and-shadows:source}{luz} por qualquer campo conta como instantânea ao lado da
do \omterm{def:g3:sound-around-us:sound}{som}: a batida vista marca o instante inicial verdadeiro com uma
precisão muito melhor do que a mão no cronômetro consegue aproveitar.
\end{solution}

\begin{solution}{exo:g7:sound-production:7}
$340 \times 1 = \qty{340}{m}$ de ida e volta: a parede está a
\qty{170}{m}. Responder \qty{340}{m} esquece que aquele segundo pagou
a ida \emph{e} a volta.
\end{solution}

\begin{solution}{exo:g7:sound-production:8}
Na água: $1500 \times 2 = \qty{3000}{m}$ de ida e volta ---
profundidade de \qty{1500}{m}. O escorregão: pegar emprestados os
\qty{340}{m/s} do \omterm{def:g2:air-around-us:air}{ar} para uma viagem submarina.
\end{solution}

\begin{solution}{exo:g7:sound-production:9}
$d = 340 \times 2.5 = \qty{850}{m}$. O clarão do fogo de artifício
deu um tiro de largada verdadeiro; o jato que não se vê não oferece
nenhum --- não há um instante de emissão para cronometrar --- e,
pior, o ronco ouvido saiu do jato segundos atrás: o ouvido aponta
para onde o jato \emph{estava}.
\end{solution}

\begin{solution}{exo:g7:sound-production:10}
A orelha submersa é servida pelo revezamento de \qty{1500}{m/s} da
água, e a orelha no deck, pelos \qty{340}{m/s} do \omterm{def:g2:air-around-us:air}{ar} --- e boa parte
do \omterm{def:g3:sound-around-us:sound}{som} da caixa nem chega a cruzar a fronteira da superfície até o
\omterm{def:g2:air-around-us:air}{ar}. \omterm{def:g7:sound-production:medium}{Meio} mais rápido, mensagem particular: o nadador vence apesar da
distância.
\end{solution}

\begin{solution}{exo:g7:sound-production:11}
Uma martelada, dois revezamentos: pelo aço e pelo \omterm{def:g2:air-around-us:air}{ar}. Aço:
$3000 \div 5000 = \qty{0.6}{s}$; \omterm{def:g2:air-around-us:air}{ar}:
$3000 \div 340 \approx \qty{8.8}{s}$ --- o clangor sai na frente do
estrondo por cerca de \qty{8.2}{s}.
\end{solution}

\begin{solution}{exo:g7:sound-production:12}
Um estalo manda o mesmo instante para os dois \omterm{def:g7:sound-production:medium}{meios}; a \qty{750}{m} o
gravador captura a chegada pela água (no hidrofone) e a chegada pelo
\omterm{def:g2:air-around-us:air}{ar} (no microfone). Tempos esperados: água
$750 \div 1500 = \qty{0.5}{s}$, \omterm{def:g2:air-around-us:air}{ar}
$750 \div 340 \approx \qty{2.2}{s}$ --- um intervalo de cerca de
\qty{1.7}{s}. Ao contrário: com $d$ e a \omterm{def:g7:motion-average-speed:formula}{velocidade} no \omterm{def:g2:air-around-us:air}{ar} conhecidos, o
intervalo medido dá o tempo de chegada pela água
($t_{\text{ar}} - \text{intervalo}$), e
$v = 750 \div 0.5 = \qty{1500}{m/s}$ --- o ritmo da água medido sem
jamais cronometrar o estalo silencioso em si.
\end{solution}

\begin{solution}{pb:g7:sound-production:1}
\textbf{1.} $340 \times 4 = \qty{1360}{m}$ de ida e volta: largura de
\qty{680}{m}.
\textbf{2.} $340 \times 1.5 = 510$; metade: \qty{255}{m}.
\textbf{3.} Primeira parede: $340 \times 1 \div 2 = \qty{170}{m}$;
segunda: $340 \times 2 \div 2 = \qty{340}{m}$; o cânion mede
$170 + 340 = \qty{510}{m}$ ao longo dessa linha.
\textbf{4.} A fórmula é exata; a mão não é: começar e parar à mão se
espalha por décimos de segundo --- dezenas de \omterm{def:g2:measuring-length:units}{metros} a
\qty{340}{m/s}. Daí a velha regra: repita, veja as leituras se
agruparem e informe o agrupamento.
\textbf{5.} $1500 \times 0.2 = \qty{300}{m}$ de ida e volta:
profundidade de \qty{150}{m}.
\textbf{6.} $1500 \times 0.4 \div 2 = \qty{300}{m}$.
\textbf{7.} Na fronteira ar--água, a maior parte de um grito é
devolvida --- pouco entra no rio, e o \omterm{ex:g4:how-sound-travels:echo}{eco} fraco do fundo precisa
cruzar a fronteira de novo para chegar a um ouvido no \omterm{def:g2:air-around-us:air}{ar}: quase nada
sobrevive. E, mesmo admitindo o \omterm{ex:g4:how-sound-travels:echo}{eco}, viagens de ida e volta de
décimos de segundo derrotam qualquer cronometragem de ouvido.
\textbf{8.} O tempo de \omterm{ex:g4:how-sound-travels:echo}{eco} se alonga bruscamente quando o barco cruza
o degrau --- fundo mais fundo, ida e volta mais longa: o sonar
desenha o penhasco debaixo d'água.
\textbf{9.} Aço: $2000 \div 5000 = \qty{0.4}{s}$; \omterm{def:g2:air-around-us:air}{ar}:
$2000 \div 340 \approx \qty{5.9}{s}$.
\textbf{10.} Os dois revezamentos começam juntos, então o intervalo
cresce em proporção à distância: cada quilômetro acrescenta um atraso
extra fixo ao revezamento lento em relação ao rápido. Meça o
intervalo, divida pela diferença por quilômetro, e a distância cai
pronta --- uma martelada só, sem precisar de cronômetro na outra
ponta.
\textbf{11.} O \omterm{def:g3:sound-around-us:sound}{som} no trilho: a neblina é uma multidão de gotinhas
que espalha a \emph{\omterm{def:g1:light-and-shadows:source}{luz}} sem esperança, mas quase não incomoda um
revezamento trancado dentro de aço \omterm{def:g3:solids-liquids-gases:solid}{sólido} (o \omterm{def:g3:sound-around-us:sound}{som} no \omterm{def:g2:air-around-us:air}{ar} também
sobrevive bem à neblina --- mas a mensagem do trilho é a mais limpa e
a mais rápida das três).
\textbf{12.} Por exemplo: ``Todo \omterm{ex:g4:how-sound-travels:echo}{eco} paga a viagem de ida e volta:
divida por dois antes de acreditar. Todo \omterm{def:g7:sound-production:medium}{meio} impõe o ritmo dele:
escolha a \omterm{def:g7:motion-average-speed:formula}{velocidade} pela multidão, nunca por hábito. E toda
cronometragem neste cânion começou com o olho --- a viagem da \omterm{def:g1:light-and-shadows:source}{luz}
conta como instantânea aqui, e é por isso que ver o martelo cair pode
servir de tiro de largada.''
\end{solution}
